\subsection{Overview}
\label{sec:method:overview}

Let an utterance be represented by three temporally aligned token sequences
$X^{t}\in\mathbb{R}^{T\times d_{t}}$ (text, last hidden states of
DeBERTa-v3-base \citep{he2023debertav3}),
$X^{a}\in\mathbb{R}^{T\times d_{a}}$ (acoustic, COVAREP
\citep{degottex2014covarep}) and
$X^{v}\in\mathbb{R}^{T\times d_{v}}$ (visual, OpenFace 2.0
\citep{baltrusaitis2018openface}).
The goal is to predict a continuous sentiment score $y\in\mathbb{R}$.

\methodname{} is built around a single design principle: the
\emph{IB-derived per-token confidence}
\(\conf=\sigmoid(-\log\sigma^{2})\in[0,1]^{T\times d_{B}}\) is a universal
control signal that can modulate fusion at three different scales
(Figure~\ref{fig:overview}):
\textbf{(i) token level} -- biasing cross-attention scores and gating values
(\S\ref{sec:method:ibca});
\textbf{(ii) sample level} -- modulating the magnitude of cross-modal injection
by the cosine similarity of the bottleneck representations
(\S\ref{sec:method:align});
\textbf{(iii) global level} -- supervising a Gumbel--softmax modality selector
through a stage-2 routing-IB loss (\S\ref{sec:method:mselector}).

\begin{figure*}[t]
\centering
\resizebox{\linewidth}{!}{%
\begin{tikzpicture}[
    node distance=4mm and 7mm,
    every node/.style={font=\small},
    box/.style={draw, rounded corners=2pt, align=center, minimum height=8mm, inner sep=3pt, fill=white},
    smallbox/.style={draw, rounded corners=2pt, align=center, minimum height=6mm, inner sep=2pt, font=\footnotesize, fill=white},
    arrow/.style={-Stealth, thick},
    dashedarrow/.style={-Stealth, thick, dashed},
    confarrow/.style={-Stealth, thick, dotted, draw=black!60},
    grouplabel/.style={font=\bfseries\footnotesize, align=center}
  ]

  % --- Inputs (left) ---
  \node[box] (xt) {Text $X^{t}$};
  \node[box, below=of xt] (xa) {Acoustic $X^{a}$};
  \node[box, below=of xa] (xv) {Visual $X^{v}$};

  % --- DeBERTa for text ---
  \node[smallbox, right=of xt] (dberta) {DeBERTa-v3};

  % --- Projections ---
  \node[smallbox, right=of dberta] (pt) {Proj$^{t}$};
  \node[smallbox, right=of xa] (pa) {Proj$^{a}$};
  \node[smallbox, right=of xv] (pv) {Proj$^{v}$};

  % --- IB encoders ---
  \node[box, right=of pt] (ibt) {IBE$^{t}$};
  \node[box, right=of pa] (iba) {IBE$^{a}$};
  \node[box, right=of pv] (ibv) {IBE$^{v}$};

  % --- Bottleneck + confidence labels ---
  \node[smallbox, right=of ibt, fill=blue!8] (bt) {$\bn^{t},\,\conf^{t}$};
  \node[smallbox, right=of iba, fill=blue!8] (ba) {$\bn^{a},\,\conf^{a}$};
  \node[smallbox, right=of ibv, fill=blue!8] (bv) {$\bn^{v},\,\conf^{v}$};

  % --- MSelector ---
  \node[box, right=12mm of ba, fill=orange!10] (msel) {MSelector\\(Gumbel-softmax)};

  % --- Routing block ---
  \node[box, right=of msel, minimum width=22mm] (route) {Soft route\\$\bn^{p},\bn^{a_{1}},\bn^{a_{2}}$};

  % --- InfoGate stack ---
  \node[box, right=of route, minimum width=30mm, minimum height=18mm, fill=green!8] (ig) {InfoGate $\times N$\\IB-guided CA\\+ alignment\\+ adaptive gate};

  % --- Head ---
  \node[box, right=of ig] (head) {Pool + LN\\+ Dropout\\+ Linear};
  \node[box, right=of head] (yhat) {$\hat{y}$};

  % --- Arrows: inputs to projections ---
  \draw[arrow] (xt) -- (dberta);
  \draw[arrow] (dberta) -- (pt);
  \draw[arrow] (xa) -- (pa);
  \draw[arrow] (xv) -- (pv);

  % --- Projections to IB encoders ---
  \draw[arrow] (pt) -- (ibt);
  \draw[arrow] (pa) -- (iba);
  \draw[arrow] (pv) -- (ibv);

  % --- IB to bottleneck/confidence ---
  \draw[arrow] (ibt) -- (bt);
  \draw[arrow] (iba) -- (ba);
  \draw[arrow] (ibv) -- (bv);

  % --- Bottlenecks into MSelector ---
  \draw[arrow] (bt.east) -- ++(4mm,0) |- (msel.west);
  \draw[arrow] (ba.east) -- (msel.west);
  \draw[arrow] (bv.east) -- ++(4mm,0) |- (msel.west);

  % --- MSelector to route ---
  \draw[arrow] (msel) -- (route);

  % --- Bottlenecks/confidence into route block (data path) ---
  \draw[arrow, dashed] (bt.east) -- ++(2mm,0) -| ([xshift=-3mm]route.north);
  \draw[arrow, dashed] (ba.east) -- ++(2mm,0) -- (route.west);
  \draw[arrow, dashed] (bv.east) -- ++(2mm,0) -| ([xshift=3mm]route.south);

  % --- Routing to InfoGate stack ---
  \draw[arrow] (route) -- (ig);

  % --- InfoGate to head ---
  \draw[arrow] (ig) -- (head);
  \draw[arrow] (head) -- (yhat);

  % --- Loss labels (below) ---
  \node[grouplabel, below=12mm of bv.south west, anchor=west, text=blue!50!black] (libtib) {$\Lvtib$, $\Lvlib$ (per-modality IB)};
  \node[grouplabel, right=15mm of libtib, text=orange!60!black] (lrib) {$\Lvrib$ (stage 2)};
  \node[grouplabel, right=8mm of lrib, text=green!50!black] (lk) {Token-/sample-/global-level\\IB modulation};
  \node[grouplabel, right=10mm of lk, text=red!60!black] (ltask) {$\Lvtask=|\hat y-y|+w_{\mathrm{mse}}(\hat y-y)^{2}$};

  % --- Connection from each IB encoder block (training-only loss) ---
  \draw[confarrow, blue!50!black] (ibt.south) to[bend right=10] (libtib.north);
  \draw[confarrow, blue!50!black] (iba.south) to[bend right=10] (libtib.north);
  \draw[confarrow, blue!50!black] (ibv.south) to[bend right=20] (libtib.north);
  \draw[confarrow, orange!60!black] (msel.south) to[bend left=10] (lrib.north);
  \draw[confarrow, green!50!black] (ig.south) -- (lk.north);
  \draw[confarrow, red!60!black] (yhat.south) to[bend left=10] (ltask.north east);

\end{tikzpicture}%
}
\caption{Overview of \methodname{}. Each modality is projected, encoded by
a variational IB encoder that exposes a per-token confidence
$\conf^{m}=\sigmoid(-\log\sigma^{2,m})$, and routed by a Gumbel-softmax
modality selector. The InfoGate stack (\S\ref{sec:method:ibca}--\ref{sec:method:gate})
uses the same confidence signal at three different scales (token, sample,
global) to modulate cross-attention, alignment-aware injection and
adaptive gating before a primary-only regression head.}
\label{fig:overview}
\end{figure*}


\subsection{Per-Modality Projection and IB Encoder}
\label{sec:method:ib}

Each modality is first projected into a shared $d_{u}$-dimensional space
with $\mathrm{Proj}^{m}(X^{m})$ realised as a linear layer plus ReLU
(text), or LayerNorm-Linear-ReLU (acoustic / visual; the LayerNorm
stabilises the very different feature statistics of the raw COVAREP /
OpenFace streams).
Each projected sequence $F^{m}\in\mathbb{R}^{T\times d_{u}}$ is then fed to
a variational IB encoder $\mathrm{IBE}^{m}$ that returns mean
$\mu^{m}$ and log-variance $\log\sigma^{2,m}$, from which a stochastic
bottleneck representation is sampled with the standard reparameterization
trick \citep{kingma2014vae,alemi2017vib}:
\begin{align}
  (\mu^{m},\log\sigma^{2,m}) &= \mathrm{IBE}^{m}(F^{m}), \\
  \bn^{m} &= \mu^{m} + \boldsymbol{\epsilon}\odot\exp(\tfrac{1}{2}\log\sigma^{2,m}),
\end{align}
with $\boldsymbol{\epsilon}\!\sim\!\mathcal{N}(0,\matr{I})$ during training and
$\bn^{m}=\mu^{m}$ at evaluation.
The IB principle \citep{tishby2000information,tishby2015dnn,alemi2017vib}
encourages the bottleneck to retain only task-relevant information,
making $\log\sigma^{2,m}$ a natural proxy for per-token uncertainty.
We therefore define a per-token, per-dimension \emph{confidence map}
\begin{equation}
  \conf^{m} \;=\; \sigmoid\!\big(-\log\sigma^{2,m}\big)\;\in\;[0,1]^{T\times d_{B}},
  \label{eq:conf}
\end{equation}
which is reused at every later modulation stage.

To make the bottleneck informative we attach an intra-modal decoder
$\mathrm{Dec}^{m}$ for each modality and minimise the
token-level IB objective
\begin{equation}
  \Lvtib \;=\; \sum_{m\in\{t,a,v\}} \KL\big(\mu^{m},\sigma^{2,m}\big) \;+\; \beta_{\mathrm{ib}}\,\mathrm{MSE}\!\big(\mathrm{Dec}^{m}(\bn^{m}),F^{m}\big),
  \label{eq:tib}
\end{equation}
where $\KL(\mu,\sigma^{2}) = -\tfrac{1}{2}\sum_{i}\big(1+\log\sigma_{i}^{2}-\mu_{i}^{2}-\sigma_{i}^{2}\big)$
is the closed-form KL against $\mathcal{N}(0,\matr{I})$, masked over valid
tokens.
A label-level IB term anchors the bottleneck to the sentiment label
through a per-modality linear head $\mathrm{LP}^{m}$:
\begin{equation}
  \Lvlib \;=\; \tfrac{1}{3}\sum_{m}\Big[\KL(\mu^{m}_{\bar{t}},\sigma^{2,m}_{\bar{t}}) + \beta_{\mathrm{ib}}\,|\mathrm{LP}^{m}(\bn^{m}_{\bar{t}})-y|\Big],
  \label{eq:lib}
\end{equation}
where $\bar{t}$ denotes mean-pooling over valid tokens.

\subsection{IB-Guided Cross-Attention}
\label{sec:method:ibca}

Cross-attention from auxiliary modality $a$ to primary modality $p$ is
the standard scaled dot-product mechanism with two confidence-driven
modifications.
For each head, queries are projected from $\bn^{p}$ and keys/values from
$\bn^{a}$ as $Q,K,V$.
Let $\bar{c}^{a}\in\mathbb{R}^{T}$ be the mean of $\conf^{a}$ over the
bottleneck dimension.
We modify the attention as
\begin{align}
  \matr{S} &= \tfrac{1}{\sqrt{d_{h}}}QK^{\top} + \gamma\,\log(\bar{c}^{a}+\varepsilon),
  \label{eq:scorebias}\\
  \matr{V}_{\!\text{eff}} &= V\,\odot\,\bar{c}^{a},
  \label{eq:valuegate}\\
  \mathrm{Attn}(Q,K,V,\conf^{a}) &= \softmax(\matr{S}+\matr{M})\matr{V}_{\!\text{eff}},
\end{align}
where $\gamma$ is a learnable scalar (initialised to $1$),
$\varepsilon\!=\!10^{-6}$ and $\matr{M}$ is the standard padding mask.
The score bias~\eqref{eq:scorebias} discourages attention from being
placed on uncertain key positions, and the value gate~\eqref{eq:valuegate}
shrinks the contribution of those positions even when the softmax
nevertheless attends to them.
Figure~\ref{fig:ibca} visualises the data path.
This design is fully back-propagatable through $\log\sigma^{2,a}$, so the
IB encoder learns to suppress its own outputs whenever they are not
useful for the downstream prediction.

\begin{figure}[t]
\centering
\resizebox{\columnwidth}{!}{%
\begin{tikzpicture}[
    node distance=3mm and 6mm,
    every node/.style={font=\footnotesize},
    box/.style={draw, rounded corners=2pt, align=center, minimum height=6mm, inner sep=2pt, fill=white},
    op/.style={draw, circle, inner sep=0pt, minimum size=4mm, fill=white},
    arrow/.style={-Stealth},
    confarrow/.style={-Stealth, dashed, draw=blue!60!black, thick}
  ]
  % primary
  \node[box] (bp) {primary\\$\bn^{p}$};
  % auxiliary
  \node[box, below=8mm of bp] (ba) {auxiliary\\$\bn^{a}$};
  % confidence
  \node[box, below=3mm of ba, fill=blue!8] (ca) {$\conf^{a}=\sigmoid(-\log\sigma^{2,a})$};

  % Q K V
  \node[box, right=of bp] (Q) {$Q=W_{Q}\bn^{p}$};
  \node[box, right=of ba] (K) {$K=W_{K}\bn^{a}$};
  \node[box, below=of K] (V) {$V=W_{V}\bn^{a}$};

  % score
  \node[box, right=12mm of Q, minimum width=24mm] (S) {$\matr{S} = \tfrac{QK^{\top}}{\sqrt{d_{h}}}\;{+}\;\gamma\log(\bar c^{a}+\varepsilon)$};

  % gate value
  \node[box, right=12mm of V, minimum width=24mm] (Vg) {$V_{\!\text{eff}} = V\odot\bar c^{a}$};

  % softmax + matmul
  \node[box, right=of S, fill=green!10] (sm) {$\softmax$};
  \node[box, right=of Vg, fill=green!10] (mm) {$\cdot$};

  % output
  \node[box, right=8mm of sm] (out) {output\\$W_{O}\,\mathrm{Attn}$};

  % --- arrows ---
  \draw[arrow] (bp) -- (Q);
  \draw[arrow] (ba) -- (K);
  \draw[arrow] (ba.east) -| ([xshift=-3mm]V.west) |- (V);

  \draw[arrow] (Q) -- (S);
  \draw[arrow] (K.east) -- ++(3mm,0) |- (S.west);

  \draw[arrow] (V) -- (Vg);

  % conf -> score bias and conf -> value gate
  \draw[confarrow] (ca.east) -| ([xshift=-3mm]S.south) -- (S.south);
  \draw[confarrow] (ca.east) -- ++(8mm,0) |- (Vg.south);

  \draw[arrow] (S) -- (sm);
  \draw[arrow] (Vg) -- (mm);
  \draw[arrow] (sm.east) -- ++(2mm,0) |- (mm.north);
  \draw[arrow] (mm.east) -- (out.west);

\end{tikzpicture}%
}
\caption{IB-Guided Cross-Attention for one head. Per-token confidence
$\conf^{a}$ from the IB encoder of the auxiliary modality is used in two
places: (i)~as a log-bias on attention scores (suppressing uncertain
keys) and (ii)~as a multiplicative gate on values
(suppressing uncertain contributions). Both effects are
back-propagatable through $\log\sigma^{2,a}$.}
\label{fig:ibca}
\end{figure}


\subsection{Alignment-Modulated Injection}
\label{sec:method:align}

Auxiliary tokens may be uninformative not only in absolute terms
(captured by $\conf^{a}$) but also \emph{relative} to the primary
modality -- e.g.\ a smiling face during a sarcastic verbal utterance.
We model this with a sample-level alignment scalar
\begin{equation}
  \alpha^{p,a}_{i} \;=\; \max\!\Big(0.3,\;\tfrac{1}{2}\big(\cos(\bar{\bn}^{p}_{i},\bar{\bn}^{a}_{i})+1\big)\Big)\in[0.3,1],
\end{equation}
where $\bar{\bn}^{m}_{i}$ is the masked-mean of $\bn^{m}$ for sample
$i$.
The lower bound at $0.3$ prevents the auxiliary path from being closed
entirely, which we found to destabilise training when the bottleneck
representations have not yet converged.
$\alpha^{p,a}_{i}$ multiplies the gated cross-attention output before it
is added to the primary stream (\S\ref{sec:method:gate}), giving a
sample-adaptive injection weight that contracts to $0.3$ when the two
modalities disagree.

\subsection{Adaptive Information Gate}
\label{sec:method:gate}

The vanilla way to mix a cross-attention output with the primary stream
is a residual addition.
We instead learn a per-sample, per-dimension gate
\begin{equation}
  \matr{g}^{a} = \sigmoid\!\big(\matr{W}_{2}\,\mathrm{ReLU}(\matr{W}_{1}[\bn^{p}\,\Vert\,\mathrm{CA}^{a}])\big),
  \label{eq:adagate}
\end{equation}
where $\mathrm{CA}^{a}$ is the IB-guided cross-attention output for
auxiliary $a$ and $[\cdot\Vert\cdot]$ denotes concatenation along the
feature axis.
The fused primary representation in one InfoGate layer is then
\begin{equation}
  \bn^{p} \leftarrow \bn^{p}+\mathrm{SA}^{p}+\sum_{a\in\{a_{1},a_{2}\}}\!\!\alpha^{p,a}\,\big(\matr{g}^{a}\odot\mathrm{CA}^{a}\big),
  \label{eq:layer_update}
\end{equation}
followed by LayerNorm and a position-wise feed-forward block.
Stacking $N$ such layers yields the InfoGate fusion stack; in all but
the last layer we additionally back-propagate primary information into
the two auxiliary streams to maintain residual alignment.

\subsection{Modality Selector and Routing-IB Loss}
\label{sec:method:mselector}

For each sample we want to choose which of the three modalities should
play the role of the primary stream.
Modality Selector (MSelector) computes a single score per modality with
adaptive attention pooling
\begin{equation}
  h^{m} \;=\; \sum_{t=1}^{T}\softmax_{t}\!\big(\matr{w}^{m\top}\bn^{m}/\sqrt{d_{B}}\big)\,\bn^{m}_{t},
\end{equation}
followed by a 2-layer MLP that produces logits
$\boldsymbol{\ell}=\mathrm{MLP}([h^{a};h^{t};h^{v}])\in\mathbb{R}^{3}$.
At training time we draw a hard one-hot vector with a Gumbel--softmax
\citep{jang2017gumbel,maddison2017concrete}
$\matr{o}=\mathrm{GumbelSoftmax}(\boldsymbol{\ell},\tau)$, kept
differentiable through the straight-through estimator, while the soft
weights $\matr{w}=\softmax(\boldsymbol{\ell}/\tau)$ scale the auxiliary
streams; $\tau$ is annealed linearly from $\tau_{\text{start}}$ to
$\tau_{\text{end}}$ over training.
At inference we use $\arg\max(\boldsymbol{\ell})$.

To prevent MSelector from collapsing onto a fixed text-centric prior
that simply mimics MAG-BERT \citep{rahman2020magbert} or
Self-MM \citep{yu2021selfmm}, we add a stage-2 routing-IB loss that
supervises the soft weights with a per-sample modality-quality target.
Specifically, we use the per-modality label-prediction errors as a
teacher signal
\begin{equation}
  \matr{q}_{i} = \softmax\!\big(-|\mathrm{LP}^{m}(\bn^{m}_{\bar{t}})-y_{i}|/\tau_{q}\big),
\end{equation}
and minimise the per-sample KL divergence on the subset of samples for
which the modalities actually differ in quality
($\mathrm{range}(\text{errors}) > 0.5$):
\begin{equation}
  \Lvrib \;=\; \mathbb{E}_{i\in\mathcal{D}}\,\KL(\matr{w}_{i}\,\Vert\,\matr{q}_{i}) \;+\; \lambda_{\!b}\,\big(\log 3 - H(\bar{\matr{w}})\big),
  \label{eq:rib}
\end{equation}
where $\bar{\matr{w}}$ is the batch-mean weight vector and the second
term is a mild entropy regulariser ($\lambda_{\!b}\!=\!0.0$ in our final
configurations).
$\Lvrib$ is only computed in stage 2 (after $\mathrm{stage1\_epochs}$
warm-up) so that the per-modality label predictors $\mathrm{LP}^{m}$ are
already meaningful before they are used as teachers.

\subsection{Training Objective}
\label{sec:method:loss}

Putting it all together, \methodname{} optimises
\begin{equation}
  \mathcal{L} \;=\; \Lvtask \;+\; \alpha_{\mathrm{ib}}\,(\Lvtib+\Lvlib) \;+\; \lambda_{\!\mathrm{rib}}\cdot \mathbb{1}_{[\text{stage}=2]}\,\Lvrib,
  \label{eq:total}
\end{equation}
with the regression task loss
\begin{equation}
  \Lvtask \;=\; |\hat y-y| + w_{\!\mathrm{mse}}\,(\hat y-y)^{2},
  \label{eq:task}
\end{equation}
where $\hat y$ is the scalar output of a LayerNorm--Dropout--Linear head
applied to the adaptively pooled primary bottleneck.
The weight $w_{\!\mathrm{mse}}$ explicitly trades MAE for Acc-2 / F1
classification quality (\S\ref{sec:exp:results}); we report two
configurations on \textsc{CMU-MOSI} that sit at different points of this
trade-off.
A two-stage schedule defines $\text{stage}=1$ if the current epoch is
below $\mathrm{stage1\_epochs}$ and $\text{stage}=2$ thereafter; only
stage 2 activates the routing supervision in
Equation~\eqref{eq:rib}.

We use AdamW \citep{loshchilov2019adamw,kingma2015adam} with two
parameter groups (a small learning rate for DeBERTa, a 10--40$\times$
larger one for the InfoGate parameters), linear warm-up over a
$\mathrm{warmup\_proportion}$ fraction of total steps, and an
exponential moving average of the model weights for evaluation after
$\mathrm{ema\_start\_epoch}$.
Hyper-parameter search uses Optuna \citep{akiba2019optuna}
(\S\ref{sec:exp:setup}).
